
\part{Background on the Hecke algebra and Soergel bimodules}

\label{pt:background}

%%%%%%%%%%%%%%%%%%%%%%%%%
\section{The Hecke Algebra}
\label{sec-hecke}
%%%%%%%%%%%%%%%%%%%%%%%%%

Background on this section can be found in \cite{Humphreys}.

%===========
\subsection{Basic definitions}
\label{subsec-basicdefinition}
%===========

Let $(W,S)$ be a Coxeter system and let $e \in W$ denote the identity. That is, $W$ is the group generated by $S$ subject to the relations:
\begin{equation} s^2=e   
\qquad \text{for all $s \in S$,}\end{equation}
\begin{equation} \ubr{sts\ldots}{m_{st}} = \ubr{tst\ldots}{m_{st}}
\label{braidreln} 
\qquad \text{for all $s \ne t \in S$.}
\end{equation}
The numbers $m_{st} = m_{ts}$ associated to each pair of simple reflections determine the group $W$, and must satisfy $m_{st} \in \{2,3,4,\ldots\}$ or $m_{st}=\infty$. When $m_{st}=\infty$, the
so-called \emph{braid relation} \eqref{braidreln} is omitted. The group $W$ is equipped with a Bruhat order $\le$ and a length function $\ell : W \to \ZM_{\ge 0}$.

For any subset $J \subset S$ the corresponding \emph{parabolic subgroup} $W_J$ is the subgroup generated by $s \in J$. Then $(W_J, J)$ is a Coxeter system with presentation induced
from that of $(W,S)$. The \emph{rank} of $W_J$ is the size of $J$. In particular, the parabolic subgroup of a single vertex is isomorphic to $S_2$, and the parabolic subgroup of a pair of
vertices is a finite or infinite dihedral group. We call $J$ \emph{finitary} if $W_J$ is finite, in which case it has a longest element $w_J$.

The Hecke algebra $\HB$ of $W$ is the free $\Zvv$-algebra generated by symbols $T_s$ for $s \in S$, modulo the following relations:

\begin{equation} T_s^2= (v^{-2}-1) T_s  + v^{-2} \label{rank1heckestd} \qquad \text{for all $s \in S$,}\end{equation}
\begin{equation} \ubr{T_sT_tT_s\ldots}{m_{st}} = \ubr{T_tT_sT_t\ldots}{m_{st}} \label{rank2heckestd} 
\qquad \text{for all $s \ne t \in S$.}
\end{equation}
It will be useful to work with a different normalization. If we set $H_s \define vT_s$ then these relations become:
\begin{equation} H_s^2= (v^{-1}-v)H_s + 1 \label{rank1heckestd'} \qquad \text{for all $s \in S$,}\end{equation}
\begin{equation} \ubr{H_sH_tH_s\ldots}{m_{st}} = \ubr{H_tH_sH_t\ldots}{m_{st}} \label{rank2heckestd'} 
\qquad \text{for all $s \ne t \in S$.}
\end{equation}

\begin{notation} 
We will use an underlined roman letter $\un{w} = (s_1,s_2, \dots, s_m)$ to denote a finite sequence of elements of $S$. Omission of the underline will denote the product $w=s_1s_2 \dots s_m$ in $W$. The \emph{length} of $\un{w}$ is $\ell(\un{w}) = m$. Note that $\ell(\un{w}) \ge \ell(w)$ with equality holding if and only if $\un{w}$ is a reduced expression for $w$. We will often abuse notation and write $\un{w} = s_1 s_2  \dots s_m$. The underline reminds us that the sequence of simple reflections, and not just their product in $W$, is important. Given $\un{w} = s_1 s_2  \dots s_m$ we write $H_{\un{w}} = H_{s_1}H_{s_2} \dots H_{s_d}$. Certainly it is possible that $H_{\un{w}} \ne H_{\un{x}}$ even though $x = w$. Later in this paper, similar notation will apply to other iterated products or tensor products.
\end{notation}

\begin{notation} The phrase ``reduced expression" dominates this paper to such an extent we have decided to shorten it to \emph{rex}. The plural of rex is rexes. \end{notation}

Given any two rexes $\un{w}$ and $\un{w}'$ for $w \in W$, it is possible to pass from $\un{w}$ to $\un{w}'$ using only braid relations. It follows from \eqref{rank2heckestd'} that the elements $H_{\un{w}}$ and $H_{\un{w}'}$ are equal, and are denoted $H_w$. We write $H_e \define 1$ for the identity of $\HB$. These elements $\{H_w\}_{w
\in W}$ form the \emph{standard basis} of $\HB$ as a $\Zvv$-module.

%===========
\subsection{The Kazhdan-Lusztig basis}
\label{subsec-KLbasis}
%===========

The Hecke algebra is equipped with a $\Z$-linear \emph{bar involution}, denoted $h \mapsto \overline{h}$ and
uniquely specified as an algebra homomorphism by $v \mapsto v^{-1}$ and $H_s \mapsto H_s^{-1}$. A simple
calculation shows that $H_s + v = H_s^{-1}+v^{-1}$, so that the element $\un{H}_s = H_s+v$ is bar-invariant.

\begin{thm}(Kazhdan-Lusztig \cite{KaL1}) There exists a unique basis $\{\un{H}_w\}_{w \in W}$ of $\HB$ as a $\Zvv$-module, called the \emph{Kazhdan-Lusztig basis}, which satisfies:
\begin{itemize}
\item $\overline{\un{H}_w}=\un{H}_w$;
\item $\un{H}_w = H_w + \sum_{x < w} h_{x,w} H_x$ where $h_{x,w} \in v\ZM[v]$.\end{itemize}
\end{thm}

The polynomials $h_{x,w} \in \ZM[v]$ are \emph{Kazhdan-Lusztig polynomials}. 

\begin{remark} In this paper we follow the normalization of Soergel \cite{Soe4} rather than the original normalization of \cite{KaL1}. In particular, the Kazhdan-Lusztig polynomials $P_{x,w}$ in \cite{KaL1} are related to the Kazhdan-Lusztig polynomials above by the formula
\[
h_{x,w} = v^{2(\ell(w) - \ell(x))}P_{x,w}(v^{-2}).
\]
\end{remark}

Given $\un{w} = s_1s_2 \dots s_m$ we set $\un{H}_{\un{w}} \define \un{H}_{s_1} \dots \un{H}_{s_m}$. Note that $\un{H}_{\un{w}} \ne \un{H}_w$ in general, but equality does hold when $\ell(\un{w}) \le 2$ and $s_1 \ne s_2$.

An element $w \in W$ is called \emph{(rationally) smooth} if $\un{H}_w = \sum_{y \le w} v^{\ell(w) -\ell(y)}H_y$. The longest element of any finite parabolic subgroup is smooth. Any
element of a rank 2 parabolic subgroup is smooth.

We pause to present three conjectures of Kazhdan and Lusztig:

\begin{conj} The polynomials $h_{x,y}$ have positive coefficients, i.e. they lie in $\ZM_{\ge 0}[v]$. \end{conj}

\begin{conj} The structure coefficients of the Kazhdan-Lusztig basis $\HB$ are positive. In other words, $\un{H}_u \un{H}_v = \sum c^w_{u,v} \un{H}_w$ for some $c^w_{u,v} \in \Nvv$. \end{conj}

These first two conjectures are commonly referred to as the \emph{Kazhdan-Lusztig positivity conjectures}. The following conjecture is the \emph{Kazhdan-Lusztig conjecture}. It and its many generalizations account to a large extent for the interest in Kazhdan-Lusztig polynomials. (For a precise statement of the Kazhdan-Lusztig conjecture see \cite{KaL1}.)

\begin{conj} (The Kazhdan-Lusztig Conjecture) If $W$ is the Weyl group of a complex semi-simple Lie algebra, then the Kazhdan-Lusztig polynomials give the multiplicities of simple modules in Verma modules in the principal block of category $\OC$.\end{conj}

These three conjectures are themselves implications of the Soergel conjecture. We will discuss this in section \ref{subsec-SoergelCatfn}.

%===========
\subsection{The presentation in the Kazhdan-Lusztig generators}
\label{subsec-KLpresentation}
%===========

As the elements $H_s$ generate $\HB$, so too do the elements $\un{H}_s$. The corresponding relations are slightly more complicated:
\begin{equation} \un{H}_s^2 = \un{H}_s (v+v^{-1}), \label{rank1hecke} \end{equation}
\begin{equation} \ubr{\un{H}_s\un{H}_t\un{H}_s\ldots}{k+1} = \sum_d  c_{k,d} \un{H}_{\ubr{sts\ldots}{d+1}} \ \mathrm{ for }\ 1 \le k+1 \le m_{st}.  \label{rank2hecke} \end{equation}
The first relation corresponds to \eqref{rank1heckestd}.
The second relation expresses a product of generators in terms of the Kazhdan-Lusztig basis within a given dihedral group. The coefficients $c_{k,d}$ appearing are decomposition
numbers for $\mathfrak{sl}_2$ tensor products. If we let $V_i$ denote the $i+1$-dimensional irreducible of $\mathfrak{sl}_2$, then $V_1^{\ot k} \bigoplus \cong V_d^{\oplus c_{k,d}}$. Roughly
speaking, this is because the Temperley-Lieb algebra at a root of unity ``categorifies" the alternating product $\un{H}_s \un{H}_t \un{H}_s \ldots$ in the dihedral Hecke algebra. More details can be
found in \cite{EDihedral}.

Relation \eqref{rank2hecke} holds even when $m_{st}=\infty$, and for any $m$ can be viewed as a definition or explicit construction of each $\un{H}_w$. When $m=\infty$, this relation does not
impose any new algebraic relations on products of $\un{H}_s$ and $\un{H}_t$. When $m < \infty$, there is one new relation on products of $\un{H}_s$ and $\un{H}_t$ coming from the fact that $\un{H}_{\ubr{sts\ldots}{m}} = \un{H}_{\ubr{tst\ldots}{m}}$.

\begin{ex} \label{rank2heckeEx} The first few examples are: \[\begin{array}{lc} m_{st}=2: & \un{H}_s\un{H}_t=\un{H}_t\un{H}_s \\ m_{st}=3: & \un{H}_s\un{H}_t\un{H}_s-\un{H}_s=\un{H}_t\un{H}_s\un{H}_t-\un{H}_t\\ m_{st}=4: & \un{H}_s\un{H}_t\un{H}_s\un{H}_t-2\un{H}_s\un{H}_t = \un{H}_t\un{H}_s\un{H}_t\un{H}_s - 2 \un{H}_t \un{H}_s \\ m_{st}=5: & \un{H}_s\un{H}_t\un{H}_s\un{H}_t\un{H}_s - 3
\un{H}_s\un{H}_t\un{H}_s + \un{H}_s = \un{H}_t\un{H}_s\un{H}_t\un{H}_s\un{H}_t - 3 \un{H}_t \un{H}_s\un{H}_t + \un{H}_t. \end{array} \] \end{ex}

%===========
\subsection{The standard trace and the defect formula}
\label{subsec-defectformula}
%===========

A \emph{trace} on $\HB$ is a $\Zvv$-linear map $\e \co \HB \to \Zvv$ satisfying $\e(hh') = \e(h'h)$ for all $h,h' \in \HB$. A straightforward calculation shows that the map $\e(\sum c_w
H_w) = c_e$ is a trace, called the \emph{standard trace}. There is a nice combinatorial formula for the standard trace of a product $\un{H}_{\un{w}}$, known
as the \emph{defect formula}, which we now discuss.

A \emph{subsequence} of $\un{w} = s_1s_2\dots s_m$ is a sequence $\pi_1 \pi_2 \dots \pi_m$ such that $\pi_i \in \{ e, s_i \}$ for all $1 \le i \le m$. Instead of working with subsequences,
we work with the equivalent datum of a sequence $\eb = \eb_1 \eb_2 \dots \eb_m$ of 1's and 0's giving the indicator function of a subsequence, which we refer to as a \emph{01-sequence}.

We can also think of $\eb$ as a roadmap for a gentle stroll through the Bruhat graph (with much pausing to admire the scenery). This \emph{Bruhat stroll} is the sequence $x_0=e, x_1,
\dots, x_m$ defined by \[ x_i \define s_1^{\eb_1} s_2^{\eb_2} \dots s_i^{\eb_i} \] for $0 \le i \le m$. We call $x_m$ the \emph{end-point} of the Bruhat stroll, and denote it by $\un{w}^{\eb}$. Alternatively, we will say that a subsequence $\eb$ of $\un{w}$ \emph{expresses} the end-point $\un{w}^{\eb}$.
The Bruhat stroll allows us to decorate each index of $\eb$ with an additional token, either U(p) or D(own). If $\eb_i=1$ so that $x_i = x_{i-1}s_i$, then we assign U to $i$ if $x_i >
x_{i-1}$ (so that we moved up in the Bruhat order at time $i$) or D if $x_i < x_{i-1}$ (so that we moved down at time $i$). If $\eb_i=0$ so that $x_i = x_{i-1}$, we glance longingly in the
direction of $x_{i-1}s_i$ but remain unmoved: we assign U or D to the index $i$ according to whether $x_{i-1}s_i > x_{i-1}$ or $x_{i-1}s_i < x_{i-1}$. The \emph{defect} of a 01-sequence
$\eb$, denoted $d(\eb)$, is defined to be the number of U0's minus the number of D0's. It measures the defect between where we longed to go and where we actually went.

\begin{ex} \label{ex:defects}
Here are some examples of subexpressions, end-points and defects:
  \begin{itemize}
  \item Suppose that $\un{w}=sss$. There are four subsequences with end-point $e$: $(U1,D1,U0)$ and $(U0,U1,D1)$ with defect $1$, $(U1,D0,D1)$ with defect $-1$, and $(U0,U0,U0)$
with defect $3$. There are four subsequences with end-point $s$: $(U1,D1,U1)$ and $(U0,U1,D0)$ with defect $1$, $(U1,D0,D0)$ with defect $-2$, and $(U0,U0,U1)$ with defect $2$.
  \item Suppose that $\un{w} = sts$ and that $m_{st}  = 3$. There are unique subexpressions with endpoints $sts, ts, st$ and $t$ with defects $0$, $1$, $1$ and $2$ respectively. There are two subexpressions $(U1, U0, D0)$ and $(U0, U0, U1)$ with end-point $s$ of defects $0$ and $2$ respectively, and two subexpressions $(U1, U0, D1)$ and $(U0, U0, U0)$ with end-point $e$ with defects $1$ and $3$ respectively.
  \end{itemize}
\end{ex}

The defect is useful because of the following lemma of Deodhar \cite{DD}:

\begin{lem} \label{lem:defects}
For any expression $\un{w}$ we have
\[
\un{H}_{\un{w}} = \sum v^{d(\eb)}H_{\un{w}^{\eb}}
\]
where the sum runs over all 01-sequences of length $\ell(\un{w})$.
\end{lem}

\begin{proof}(Sketch) It is an straightforward consequence of \eqref{rank1heckestd'} that in $\HB$ we have the relation
\[
H_x\un{H}_s = \begin{cases} H_{xs} + vH_{x} & \text{if $xs > x$,} \\
H_{xs} + v^{-1}H_x & \text{if $xs < x$.} \end{cases}
\]
We conclude that if the lemma is true for $\un{w} = \un{x}$ it is also true for $\un{w} = \un{x}s$. The result now follows by induction.
\end{proof}

We now come to the defect formula for the trace:

\begin{cor}\label{cor:defects}
For any expression $\un{w}$ we have
\[ \e(\un{H}_{\un{w}}) = \sum v^{d(\eb)} \]
where the sum is over all 01-sequences expressing the identity element.\end{cor}

\begin{ex} We continue Example \ref{ex:defects} and check Lemma \ref{lem:defects}:
  \begin{itemize}
  \item Using relation \eqref{rank1hecke} we see that
\[
\un{H}_s^3=(v+v^{-1})^2\un{H}_s = (v^{-2} + 2 + v^2)H_s + (v^{-1}+2v+v^3)H_e.
\]
  \item Using direct calculation, or smoothness and \eqref{rank2hecke}, or Lemma \ref{lem:defects} we obtain
\[
\un{H}_s\un{H}_t\un{H}_s = H_{sts} + vH_{ts} + vH_{st} + v^2H_t + (1+v^2)H_s + (v+v^3)H_e.
\]
  \end{itemize}
\end{ex}

It will be important later that the set of subsequences $\eb$ of a fixed expression $\un{w}$ is equipped with a partial order, the \emph{path dominance order}. Let $\eb$ and $\fb$ be two 01-sequences and let their corresponding Bruhat strolls be $x_0, x_1, \dots, x_m$ and $y_0, y_1, \dots, y_m$. We say that $\eb \ge \fb$ if $x_i \ge y_i$ for all $0 \le i \le m$. Clearly if $\eb \ge \fb$ then the end-point of $\eb$ is greater than the end-point of $\fb$. The path dominance order restricts to a partial order on the set of subsequences with fixed end-point.

Let $\w$ be the $\Z$-linear antiinvolution for which $\w(\un{H}_s)=\un{H}_s$ and $\w(v)=v^{-1}$. The standard trace gives rise to the \emph{standard pairing} $\HC \times \HC \to \Zvv$,
defined by $(a,b) = \e(b\w(a))$. This pairing is \emph{semilinear} over $\Zvv$; that is $(v^{-1}a,b)=(a,vb)=v(a,b)$ for all $a, b \in \HB$. Under this pairing, $b_s$ is self-biadjoint, i.e. \[(b_s x,y) = (x,b_s y),\qquad \qquad (x b_s,y) = (x,y b_s). \]

\begin{remark} \label{rmk:traces pairings} The formula $(a,b) = \e(b\w(a))$ can be used both ways, to define a pairing from a trace or vice versa. One can see that the $\Zvv$-module of all
semilinear pairings with self-biadjoint $b_s$ is isomorphic to the module of all $\Zvv$-linear traces. Any such pairing is determined by the values $\e(b_\xb)=(1,b_\xb)$ over all sequences
$\xb$. \end{remark}
